\begingroup
\renewcommand{\arraystretch}{1.45}
\fancyhead[L]{\small ARITHMETIC–GEOMETRIC MEAN.}
\Needspace{28\baselineskip}
\AGMArticle{19}

[The relations between the differentials of the quantities under investigation are particularly simple when $a$ and $b$ remain unchanged; it follows immediately from the equations expressing the conditions on terms with the same index that the expression
$\frac{1}{2^n}\frac{1}{b^{\mathrm n}}\surd{\frac{\beta^{\mathrm n}\delta_{\mathrm n}}{\alpha^{\mathrm n}\gamma^{\mathrm n}}}\cdot\mathrm{d}\log\frac{\delta^{\mathrm n}}{\beta^{\mathrm n}}$
does not change its value when $\alpha,\allowbreak \beta,\allowbreak \gamma,\allowbreak \delta$ are successively substituted for $\beta,\allowbreak \gamma,\allowbreak \delta,\allowbreak \alpha$ in it. The two expressions thus obtained can be combined into an expression for the index $n+1$ corresponding to the first one; the value of this expression is therefore also independent of the index $n$. Passing to the limit and comparing with the differentials of the quotients formed in various ways gives
\Needspace{24\baselineskip}
\[
\begin{aligned}\dfrac{1}{k}\mathrm{d}\log\eta
&=\dfrac{1}{a}\surd{\dfrac{\alpha\delta}{\beta\gamma}}\cdot\mathrm{d}\log\dfrac{\delta}{\alpha}
\\[7pt]&=\dfrac{1}{b}\surd{\dfrac{\beta\delta}{\alpha\gamma}}\cdot\mathrm{d}\log\dfrac{\delta}{\beta}
\\[7pt]&=\dfrac{1}{c}\surd{\dfrac{\gamma\delta}{\alpha\beta}}\cdot\mathrm{d}\log\dfrac{\delta}{\gamma}
\\[7pt]&=\dfrac{1}{c}\surd{\dfrac{\delta'}{\beta'}}\cdot\mathrm{d}\log\dfrac{\delta\delta}{\delta'}
\\[7pt]&=\dfrac{1}{a}\surd{\dfrac{\beta\gamma}{\alpha\delta}}\cdot\mathrm{d}\log\dfrac{\gamma}{\beta}
\\[7pt]&=\dfrac{1}{b}\surd{\dfrac{\alpha\gamma}{\beta\delta}}\cdot\mathrm{d}\log\dfrac{\gamma}{\alpha}
\\[7pt]&=\dfrac{1}{c}\surd{\dfrac{\alpha\beta}{\gamma\delta}}\cdot\mathrm{d}\log\dfrac{\alpha}{\beta}
\\[7pt]&=\dfrac{1}{c}\surd{\dfrac{\beta'}{\delta'}}\cdot\mathrm{d}\log\dfrac{\beta'}{\beta\beta}\end{aligned}
\]
\Needspace{16\baselineskip}
a relation that therefore always holds when the indices of $a,\allowbreak b,\allowbreak c,\allowbreak \alpha,\allowbreak \beta,\allowbreak \gamma,\allowbreak \delta,\allowbreak \eta$ are increased by the same number of units. From this follows the theorem recorded by \textsc{Gauss}:]
\Needspace{12\baselineskip}
\[
\begin{aligned}\int\dfrac{\mathrm{d}U}{\surd{(aa\sin U^2+bb\cos U^2)}}
&=\int\dfrac{\mathrm{d}V}{\surd{(aa\cos V^2+bb\sin V^2)}}\\[6pt]&=\dfrac{u}{k}\end{aligned}
\]

\Needspace{22\baselineskip}
\SourceLeaf{400}
[With the aid of the equation for $\mathrm{d}\log\frac{\beta'}{\beta\beta}$, repeated application of it immediately gives the series for the differentials of the individual quantities $\alpha,\allowbreak \beta,\allowbreak \gamma,\allowbreak \delta$; for a negative $\delta$, the result is:
\Needspace{22\baselineskip}
\[
\begin{aligned}\dfrac{k}{2}\dfrac{\mathrm{d}\log\dfrac{\alpha}{x}}{\mathrm{d}u}+2c'\sin 2V'
&=\dfrac{k}{2}\dfrac{\mathrm{d}\log\dfrac{\beta}{x}}{\mathrm{d}u}
\\[7pt]&=\dfrac{k}{2}\dfrac{\mathrm{d}\log\dfrac{\gamma}{x}}{\mathrm{d}u}+a\dfrac{\sin V}{\cos U}
\\[7pt]&=\dfrac{k}{2}\dfrac{\mathrm{d}\log\dfrac{\delta}{x}}{\mathrm{d}u}-b\dfrac{\cos V}{\sin U}
\\[7pt]&=c'\sin 2V'+c''\sin 2^2V''+c'''\sin 2^3V'''+.\;.\;.\text{]}\end{aligned}
\]

\Needspace{24\baselineskip}
\AGMArticle{20}

[The differentials of the second order, which refer to the quantity $\eta$ as the sole independent variable, can be derived immediately by differentiating the differential equation of the first order, and the determination of their values can be represented in such a way that expressions of the form
\Needspace{12\baselineskip}
\[
\dfrac{\mathrm{d}}{\mathrm{d}\log\eta}
\left[\dfrac{\mathrm{d}\log\dfrac{\alpha}{\beta}}{\mathrm{d}\log\eta}\right]
+\tfrac12\dfrac{\mathrm{d}\log\dfrac{\alpha\beta}{\gamma\delta}}{\mathrm{d}\log\eta}
\cdot\dfrac{\mathrm{d}\log\dfrac{\alpha}{\beta}}{\mathrm{d}\log\eta}
\]
vanish for $\alpha,\allowbreak \beta,\allowbreak \gamma,\allowbreak \delta$ bearing any common index, in whatever order $\alpha,\allowbreak \beta,\allowbreak \gamma,\allowbreak \delta$ may be substituted in them. Multiplying the numerator and denominator inside the logarithm to be differentiated twice by $\beta$, replacing $\alpha\beta$ by $b'\beta'$ and the derivatives of the first order by their values gives, for two successive indices:
\Needspace{12\baselineskip}
\[
\dfrac{\mathrm{dd}\log\dfrac{\beta'}{\beta\beta}}{(\mathrm{d}\log\eta)^2}
+2\dfrac{a'c'}{kk}\dfrac{\delta'}{\beta'}
-2\dfrac{ac}{kk}\dfrac{\delta}{\beta}-\tfrac12\dfrac{cc}{kk}=0
\]
\Needspace{24\baselineskip}
and from this, by repeated application of this equation and with the aid of the series found in Article 13 for $\mathrm{d}\log\mathrm{M}(a,\allowbreak b)$, the equation
\Needspace{22\baselineskip}
\[
\begin{aligned}\dfrac{\mathrm{d}\log\dfrac{b}{k}}{\mathrm{d}\log y}
&=\dfrac{\mathrm{dd}\log\dfrac{\alpha}{x}}{(\mathrm{d}\log\eta)^2}-\tfrac12\dfrac{ac}{kk}\dfrac{\gamma}{\alpha}
\\[7pt]&=\dfrac{\mathrm{dd}\log\dfrac{\beta}{x}}{(\mathrm{d}\log\eta)^2}+\tfrac12\dfrac{ac}{kk}\dfrac{\delta}{\beta}
\\[7pt]&=\dfrac{\mathrm{dd}\log\dfrac{\gamma}{x}}{(\mathrm{d}\log\eta)^2}-\tfrac12\dfrac{ac}{kk}\dfrac{\alpha}{\gamma}
\\[7pt]&=\dfrac{\mathrm{dd}\log\dfrac{\delta}{x}}{(\mathrm{d}\log\eta)^2}+\tfrac12\dfrac{ac}{kk}\dfrac{\beta}{\delta}\end{aligned}
\]
which remains valid when $a,\allowbreak b,\allowbreak c,\allowbreak \alpha,\allowbreak \beta,\allowbreak \gamma,\allowbreak \delta,\allowbreak x,\allowbreak y,\allowbreak \eta$ are given any common index different from zero.

\Needspace{16\baselineskip}
\SourceLeaf{401}
Combining this result with the expansion of $\mathrm{d}\log\frac{\alpha}{x}$ found in the preceding article gives the evaluation of the integral found by \textsc{Gauss}:]
\Needspace{16\baselineskip}
\[
\begin{aligned}
\int\surd{(aa\sin U^2+bb\cos U^2)}\,\mathrm{d}U
\\[7pt]&=\int\dfrac{aabb\,\mathrm{d}V}{(aa\cos V^2+bb\sin V^2)^{\frac32}}\\
&=\dfrac{u}{k}\{a'a'-2c''c''-4c'''c'''-8c''''c''''-.\;.\}\\
&\quad+c'\sin 2V'-c''\sin 4V''\\&\quad-c'''\sin 8V'''-c''''\sin 16V''''-.\;.\;.\end{aligned}
\]

\clearpage
\endgroup
